\maketitle
\makesignature

\ifproject
\begin{abstractTH}
\sloppy

% เขียนบทคัดย่อของโครงงานที่นี่
\hspace{\parindent}ปัจจุบันมหาวิทยาลัยเชียงใหม่มีการติดตั้งระบบโซลาร์เซลล์แล้วแต่ยังพบความไม่สอดคล้องระหว่างการใช้ไฟฟ้าและปริมาณการผลิตไฟฟ้าจากเซลล์แสงอาทิตย์ (Solar Cell) 
เนื่องจากพลังงานที่ผลิตได้ขึ้นอยู่กับสภาพอากาศมีการเปลี่ยนแปลงตลอดเวลา
อย่างไรก็ตาม ระบบในปัจจุบันยังไม่มีข้อมูลสภาพอากาศที่ละเอียดเพียงพอสำหรับการวิเคราะห์หรือคาดการณ์การผลิตพลังงาน ส่งผลให้ไม่สามารถวางแผนการใช้พลังงานได้อย่างมีประสิทธิภาพ

เพื่อแก้ไขปัญหานี้ โครงการจึงพัฒนาสถานีตรวจวัดสภาพอากาศขนาดย่อมสำหรับเก็บข้อมูลสิ่งแวดล้อมที่มีผลต่อการผลิตไฟฟ้าจากเซลล์แสงอาทิตย์ (Solar Cell) โดยใช้เซ็นเซอร์วัดอุณหภูมิ ความชื้นสัมพัทธ์ ความกดอากาศ ความเข้มของแสง อุณภูมิแผงโซลาร์เซลล์ ความเร็วลม และทิศทางลม 
ข้อมูลที่ได้จะถูกนำไปใช้ในการวิเคราะห์และสร้างแบบจำลองเพื่อทำนายการผลิตพลังงานจากเซลล์แสงอาทิตย์ (Solar Cell) ล่วงหน้า 24 ชั่วโมง ซึ่งจะช่วยให้สามารถวางแผนการใช้งานพลังงานและบริหารจัดการพลังงานได้อย่างมีประสิทธิภาพมากขึ้นซึ่งไม่เพียงแต่ช่วยลดค่าใช้จ่ายด้านพลังงาน แต่ยังเป็นการสนับสนุนนโยบายการลดการปล่อยก๊าซเรือนกระจกและมุ่งสู่ความเป็นกลางทางคาร์บอนในอาคารของมหาวิทยาลัย
\end{abstractTH}

\begin{abstract}
    \hspace{\parindent}Currently, Chiang Mai University has implemented solar cell systems to promote renewable energy. However, there is still a discrepancy between electricity consumption and the actual power generated from solar energy, primarily due to the intermittent nature of solar power which fluctuates according to weather conditions. Furthermore, existing systems lack high-resolution meteorological data necessary for accurate analysis or forecasting, leading to challenges in efficient energy planning.
 
    To address this issue, the project developed a compact weather station to collect environmental data affecting solar power generation, including temperature, relative humidity, air pressure, solar intensity, solar panel temperature, wind speed, and wind direction. The gathered data is utilized for analysis and modeling to predict solar energy output 24 hours in advance. This enables more efficient energy planning and management, which not only reduces energy costs but also supports greenhouse gas reduction policies and the university's goal of achieving carbon neutrality.
\end{abstract}

\iffalse
\begin{dedication}
This document is dedicated to all Chiang Mai University students.

Dedication page is optional.
\end{dedication}
\fi % \iffalse

\begin{acknowledgments}
\sloppy
\hspace{\parindent}โครงการ ระบบบริหารการจัดการพลังงานอัจฉริยะเพื่อมุ่งสู่ความเป็นกลางทางคาร์บอนภายในอาคาร (Smart Energy Management System for Carbon Neutrality in Buildings) ไม่อาจสําเร็จลุล่วงได้ หากปราศจากแรงสนับสนุน คําแนะนํา และความทุ่มเทจาก
หลายฝ่ายทั้ง ผศ.ดร.ภาสกร แช่มประเสริฐ อาจารย์ที่ปรึกษาที่ได้สละเวลาอันมีค่าในการให้คำปรึกษา แนะนำแนวคิด ตลอดจนถ่ายทอดความรู้และประสบการณ์ที่เป็นประโยชน์ยิ่ง ทำให้คณะผู้จัดทำสามารถแก้ไขปัญหาและดำเนินงานได้อย่างมีประสิทธิภาพ

ขอขอบพระคุณ ผศ.ดร.พฤษภ์ บุญมา และ ผศ.ดร.ยุทธพงษ์ สมจิต 
อาจารย์กรรมการสอบโครงการ ที่สละเวลาในวันเวลาสอบเพื่อตรวจสอบ ให้ข้อเสนอแนะ และแนะนําแนวทางในการ
ปรับปรุงระบบให้มีคุณภาพและความสมบูรณ์ยิ่งขึ้น

นอกจากนี้ คณะผู้จัดทำขอขอบคุณรุ่นพี่และเพื่อนที่คอยช่วยทดสอบระบบ และคอยให้คําแนะนําต่างๆ
คณะผู้จัดทำขอถือโอกาสนี้แสดงความขอบคุณจากใจจริงต่อทุกภาคส่วนที่มีส่วนร่วมในการผลักดันให้โครงการนี้เกิดขึ้น

\acksign{2026}{3}{20}
\end{acknowledgments}%
\fi % \ifproject

\contentspage

\ifproject
\figurelistpage

\tablelistpage
\fi % \ifproject

% \abbrlist % this page is optional

% \symlist % this page is optional

% \preface % this section is optional
